% from ega4-21c-fr.tex
\oldpage[IV-4]{326}
Comme les hypothèses et la conclusion sont locales sur $S$ et sur $X$ en vertu de
(21.13.6, (i)), on peut se borner au cas où $S$ et $X$ sont affines, $f$ étant donc de
présentation finie. Puisque les fibres de $f$ sont régulières, l'hypothèse sur $Y$
entraîne que, pour tout point $x\in Y$, on a
\[
  \operatorname{prof}(\mathcal O_{X_{f(x)},x})
  =\dim(\mathcal O_{X_{f(x)},x})\geq1;
\]
il résulte de (19.9.8) que l'homomorphisme canonique
$\mathcal O_X\to j_*(\mathcal O_U)$ est \emph{injectif}. De plus, remplaçant $Y$ par
une partie fermée plus grande suivant la méthode de la démonstration de (19.9.8), et
utilisant (21.13.6, (ii)), on peut supposer que $Y$ est défini par un idéal \emph{de
type fini} de l'anneau de $X$, donc l'ouvert $U=X-Y$ est \emph{quasi-compact et
quasi-séparé} et l'ensemble fermé $Y$ est \emph{constructible}.

Pour prouver (ii), il suffit, pour tout $\mathcal O_U$-Module inversible $\mathscr L$
et tout point $x\in Y$ donnés, d'établir l'existence d'un voisinage ouvert $V$ de $x$
dans $X$ tel que : $1^{\mathrm o}$ l'homomorphisme canonique
$\mathcal O_V\to(j_V)_*(\mathcal O_{U\cap V})$ est surjectif;
$2^{\mathrm o}$ il existe un $\mathcal O_V$-Module inversible $\mathscr L'_V$ tel que
$\mathscr L'_V|(U\cap V)=\mathscr L|(U\cap V)$ ((21.13.5) et (21.13.3)).

Posons $s=f(x)$; nous allons voir qu'on peut se borner au cas où
$S=S_0=\operatorname{Spec}(\mathcal O_{S,s})$. Soit en effet $(S_\nu)$ un système
fondamental de voisinages ouverts affines de $s$ dans $S$, et posons
$X_\nu=f^{-1}(S_\nu)$, $Y_\nu=Y\cap X_\nu$, $U_\nu=X_\nu-Y_\nu$,
$j_\nu:U_\nu\to X_\nu$ étant l'injection canonique; $X_0=X\times_S S_0$ est alors
limite projective des $X_\nu$ (8.1.2, $a$)); supposons la proposition vraie pour le
morphisme $f_0:X_0\to S_0$ et pour $Y_0=Y\cap X_0$, et posons $U_0=X_0-Y_0$,
$j_0:U_0\to X_0$ étant l'injection canonique. La projection
$p_\nu:X_0\to X_\nu$ étant un morphisme affine, on a
\[
  (j_0)_*(\mathcal O_{U_0})
  =p_\nu^*((j_\nu)_*(\mathcal O_{U_\nu}))
  \qquad(\mathbf{II}, 1.5.2)
\]
et l'homomorphisme canonique
$\rho_0:\mathcal O_{X_0}\to(j_0)_*(\mathcal O_{U_0})$ n'est autre que
$p_\nu^*(\rho_\nu)$, où
$\rho_\nu:\mathcal O_{X_\nu}\to(j_\nu)_*(\mathcal O_{U_\nu})$ est l'homomorphisme
canonique; comme par hypothèse $\rho_0$ est surjectif, il en est de même de $\rho_\nu$
pour $\nu$ assez grand (8.5.7). Soit d'autre part $\mathscr L_0$ le
$\mathcal O_{U_0}$-Module inversible restriction de $\mathscr L$, et notons que les
$X_\nu$ sont affines, donc quasi-compacts et quasi-séparés; puisque par hypothèse il
existe un $\mathcal O_{X_0}$-Module inversible $\mathscr L'_0$ tel que
$\mathscr L'_0|U_0=\mathscr L_0$, il existe un $\nu$ assez grand et un
$\mathcal O_{X_\nu}$-Module quasi-cohérent $\mathscr L'_\nu$ de présentation finie
tels que $\mathscr L'_0$ soit isomorphe à $p_\nu^*(\mathscr L'_\nu)$ (8.5.2, (ii));
de plus (8.5.5) on peut supposer que $\mathscr L'_\nu$ est inversible. Enfin, comme
les $U_\nu$ sont quasi-compacts et quasi-séparés et que
$\mathscr L'_0|U_0=\mathscr L_0$, il existe $\nu$ assez grand tel que
$\mathscr L'_\nu|U_\nu$ et $\mathscr L|U_\nu$ soient isomorphes (8.5.2.5).

